% Options for packages loaded elsewhere
% Options for packages loaded elsewhere
\PassOptionsToPackage{unicode}{hyperref}
\PassOptionsToPackage{hyphens}{url}
%
\documentclass[
  11pt,
  letterpaper,
  DIV=11,
  numbers=noendperiod,
  twoside]{scrartcl}
\usepackage{xcolor}
\usepackage[left=1.1in, right=1in, top=0.8in, bottom=0.8in,
paperheight=9.5in, paperwidth=7in, includemp=TRUE, marginparwidth=0in,
marginparsep=0in]{geometry}
\usepackage{amsmath,amssymb}
\setcounter{secnumdepth}{3}
\usepackage{iftex}
\ifPDFTeX
  \usepackage[T1]{fontenc}
  \usepackage[utf8]{inputenc}
  \usepackage{textcomp} % provide euro and other symbols
\else % if luatex or xetex
  \usepackage{unicode-math} % this also loads fontspec
  \defaultfontfeatures{Scale=MatchLowercase}
  \defaultfontfeatures[\rmfamily]{Ligatures=TeX,Scale=1}
\fi
\usepackage{lmodern}
\ifPDFTeX\else
  % xetex/luatex font selection
  \setmathfont[]{Garamond-Math}
\fi
% Use upquote if available, for straight quotes in verbatim environments
\IfFileExists{upquote.sty}{\usepackage{upquote}}{}
\IfFileExists{microtype.sty}{% use microtype if available
  \usepackage[]{microtype}
  \UseMicrotypeSet[protrusion]{basicmath} % disable protrusion for tt fonts
}{}
\usepackage{setspace}
% Make \paragraph and \subparagraph free-standing
\makeatletter
\ifx\paragraph\undefined\else
  \let\oldparagraph\paragraph
  \renewcommand{\paragraph}{
    \@ifstar
      \xxxParagraphStar
      \xxxParagraphNoStar
  }
  \newcommand{\xxxParagraphStar}[1]{\oldparagraph*{#1}\mbox{}}
  \newcommand{\xxxParagraphNoStar}[1]{\oldparagraph{#1}\mbox{}}
\fi
\ifx\subparagraph\undefined\else
  \let\oldsubparagraph\subparagraph
  \renewcommand{\subparagraph}{
    \@ifstar
      \xxxSubParagraphStar
      \xxxSubParagraphNoStar
  }
  \newcommand{\xxxSubParagraphStar}[1]{\oldsubparagraph*{#1}\mbox{}}
  \newcommand{\xxxSubParagraphNoStar}[1]{\oldsubparagraph{#1}\mbox{}}
\fi
\makeatother


\usepackage{longtable,booktabs,array}
\newcounter{none} % for unnumbered tables
\usepackage{calc} % for calculating minipage widths
% Correct order of tables after \paragraph or \subparagraph
\usepackage{etoolbox}
\makeatletter
\patchcmd\longtable{\par}{\if@noskipsec\mbox{}\fi\par}{}{}
\makeatother
% Allow footnotes in longtable head/foot
\IfFileExists{footnotehyper.sty}{\usepackage{footnotehyper}}{\usepackage{footnote}}
\makesavenoteenv{longtable}
\usepackage{graphicx}
\makeatletter
\newsavebox\pandoc@box
\newcommand*\pandocbounded[1]{% scales image to fit in text height/width
  \sbox\pandoc@box{#1}%
  \Gscale@div\@tempa{\textheight}{\dimexpr\ht\pandoc@box+\dp\pandoc@box\relax}%
  \Gscale@div\@tempb{\linewidth}{\wd\pandoc@box}%
  \ifdim\@tempb\p@<\@tempa\p@\let\@tempa\@tempb\fi% select the smaller of both
  \ifdim\@tempa\p@<\p@\scalebox{\@tempa}{\usebox\pandoc@box}%
  \else\usebox{\pandoc@box}%
  \fi%
}
% Set default figure placement to htbp
\def\fps@figure{htbp}
\makeatother


% definitions for citeproc citations
\NewDocumentCommand\citeproctext{}{}
\NewDocumentCommand\citeproc{mm}{%
  \begingroup\def\citeproctext{#2}\cite{#1}\endgroup}
\makeatletter
 % allow citations to break across lines
 \let\@cite@ofmt\@firstofone
 % avoid brackets around text for \cite:
 \def\@biblabel#1{}
 \def\@cite#1#2{{#1\if@tempswa , #2\fi}}
\makeatother
\newlength{\cslhangindent}
\setlength{\cslhangindent}{1.5em}
\newlength{\csllabelwidth}
\setlength{\csllabelwidth}{3em}
\newenvironment{CSLReferences}[2] % #1 hanging-indent, #2 entry-spacing
 {\begin{list}{}{%
  \setlength{\itemindent}{0pt}
  \setlength{\leftmargin}{0pt}
  \setlength{\parsep}{0pt}
  % turn on hanging indent if param 1 is 1
  \ifodd #1
   \setlength{\leftmargin}{\cslhangindent}
   \setlength{\itemindent}{-1\cslhangindent}
  \fi
  % set entry spacing
  \setlength{\itemsep}{#2\baselineskip}}}
 {\end{list}}
\usepackage{calc}
\newcommand{\CSLBlock}[1]{\hfill\break\parbox[t]{\linewidth}{\strut\ignorespaces#1\strut}}
\newcommand{\CSLLeftMargin}[1]{\parbox[t]{\csllabelwidth}{\strut#1\strut}}
\newcommand{\CSLRightInline}[1]{\parbox[t]{\linewidth - \csllabelwidth}{\strut#1\strut}}
\newcommand{\CSLIndent}[1]{\hspace{\cslhangindent}#1}



\setlength{\emergencystretch}{3em} % prevent overfull lines

\providecommand{\tightlist}{%
  \setlength{\itemsep}{0pt}\setlength{\parskip}{0pt}}



 


\setlength\heavyrulewidth{0ex}
\setlength\lightrulewidth{0ex}
\usepackage[automark]{scrlayer-scrpage}
\clearpairofpagestyles
\cehead{
  Claude
  }
\cohead{
  Fairness Has Memory
  }
\ohead{\bfseries \pagemark}
\cfoot{}
\makeatletter
\newcommand*\NoIndentAfterEnv[1]{%
  \AfterEndEnvironment{#1}{\par\@afterindentfalse\@afterheading}}
\makeatother
\NoIndentAfterEnv{itemize}
\NoIndentAfterEnv{enumerate}
\NoIndentAfterEnv{description}
\NoIndentAfterEnv{quote}
\NoIndentAfterEnv{equation}
\NoIndentAfterEnv{longtable}
\NoIndentAfterEnv{abstract}
\renewenvironment{abstract}
 {\vspace{-1.25cm}
 \quotation\small\noindent\emph{Abstract}:}
 {\endquotation}
\newfontfamily\tfont{EB Garamond}
\addtokomafont{disposition}{\rmfamily}
\addtokomafont{descriptionlabel}{\rmfamily}
\addtokomafont{title}{\normalfont\itshape}
\let\footnoterule\relax

\makeatletter
\renewcommand{\@maketitle}{%
  \newpage
  \null
  \vskip 2em%
  \begin{center}%
  \let \footnote \thanks
    {\itshape\huge\@title \par}%
    \vskip 0.5em%  % Reduced from default
    {\large
      \lineskip 0.3em%  % Reduced from default 0.5em
      \begin{tabular}[t]{c}%
        \@author
      \end{tabular}\par}%
    \vskip 0.5em%  % Reduced from default
    {\large \@date}%
  \end{center}%
  \par
  }
\makeatother
\RequirePackage{lettrine}

\renewenvironment{abstract}
 {\quotation\small\noindent\emph{Abstract}:}
 {\endquotation\vspace{-0.02cm}}

\setmainfont{EB Garamond Math}[
  BoldFont = {EB Garamond SemiBold},
  ItalicFont = {EB Garamond Italic},
  RawFeature = {+smcp},
]

\newfontfamily\scfont{EB Garamond Regular}[RawFeature=+smcp]
\renewcommand{\textsc}[1]{{\scfont #1}}

\renewcommand{\LettrineTextFont}{\scfont}
% Number theorem-like environments straight through (Theorem 1, Theorem 2, ...)
% rather than by section (1.1, 1.2, ...), which is Quarto's LaTeX default.
% Guarded, since Quarto only defines a counter for environments the document uses.
\makeatletter
\AtBeginDocument{%
  \@for\bw@thm:={theorem,lemma,corollary,proposition,conjecture,definition,%
                 example,exercise,refremark,refsolution}\do{%
    \@ifundefined{c@\bw@thm}{}{\expandafter\counterwithout\expandafter{\bw@thm}{section}}%
  }%
}
\makeatother
\KOMAoption{captions}{tableheading}
\makeatletter
\@ifpackageloaded{caption}{}{\usepackage{caption}}
\AtBeginDocument{%
\ifdefined\contentsname
  \renewcommand*\contentsname{Table of contents}
\else
  \newcommand\contentsname{Table of contents}
\fi
\ifdefined\listfigurename
  \renewcommand*\listfigurename{List of Figures}
\else
  \newcommand\listfigurename{List of Figures}
\fi
\ifdefined\listtablename
  \renewcommand*\listtablename{List of Tables}
\else
  \newcommand\listtablename{List of Tables}
\fi
\ifdefined\figurename
  \renewcommand*\figurename{Figure}
\else
  \newcommand\figurename{Figure}
\fi
\ifdefined\tablename
  \renewcommand*\tablename{Table}
\else
  \newcommand\tablename{Table}
\fi
}
\@ifpackageloaded{float}{}{\usepackage{float}}
\floatstyle{ruled}
\@ifundefined{c@chapter}{\newfloat{codelisting}{h}{lop}}{\newfloat{codelisting}{h}{lop}[chapter]}
\floatname{codelisting}{Listing}
\newcommand*\listoflistings{\listof{codelisting}{List of Listings}}
\makeatother
\makeatletter
\makeatother
\makeatletter
\@ifpackageloaded{caption}{}{\usepackage{caption}}
\@ifpackageloaded{subcaption}{}{\usepackage{subcaption}}
\makeatother
\usepackage{bookmark}
\IfFileExists{xurl.sty}{\usepackage{xurl}}{} % add URL line breaks if available
\urlstyle{same}
% fallback for those not using the hyperref driver hyperxmp:
\makeatletter
\@ifundefined{xmpquote}{\newcommand{\xmpquote}[1]{#1}}{}
\makeatother
\hypersetup{
  pdftitle={Fairness Has Memory},
  hidelinks,
  pdfcreator={LaTeX via pandoc}}


\title{Fairness Has Memory}
\author{}
\date{2026}
\begin{document}
\maketitle
\begin{abstract}
When a scarce, indivisible good is allocated by lottery among people
with equal claims, the standard theory of fairness says the lottery is
the fairest thing available. When the same good is allocated the same
way year after year to largely the same people, the theory has nothing
further to say: each lottery is fair, so the sequence is fair. I argue
that this is a mistake. A sequence of individually fair lotteries can be
an unfair allocation, and where claimants stay in the pool after being
served, fairness requires that their chances in later rounds depend on
the history of their unsatisfied claims. The argument uses only premises
the standard theory already accepts: that a divisible good must be
divided in proportion to claims, and that a lottery is a second-best
surrogate for division. Over a sequence, the good becomes divisible. I
then show what memory should track, in units of the good where later
units are as good as earlier ones and in waiting where they are not, and
how heavily it may press: only as hard as leaves no claimant's expected
share, over whatever stay she turns out to have, below what her claim
warrants. Where the good recurs, rules that forget and rules that
remember lexically are unfair in opposite directions; where winners
leave the pool and some will leave unserved, the lottery that forgets is
what fairness requires.
\end{abstract}


\setstretch{1.1}
\section{Two lotteries}\label{sec-intro}

Every year Nevada and New Mexico each have more people who want to hunt
bighorn sheep than they have sheep to spare, so each holds a lottery.
Nevada remembers: an unsuccessful applicant earns a bonus point, her
points are squared to fix how many random numbers are drawn for her the
next year, and her lowest number is her entry, so ten years of losing
gives her a hundred and one draws against a newcomer's one.\footnote{Nevada
  Department of Wildlife, ``Bonus Point Program,''
  https://www.ndow.org/blog/bonus-point-program/.} New Mexico forgets:
``New Mexico does NOT grant preference to applicants who were
unsuccessful in previous drawings.''\footnote{New Mexico Department of
  Game and Fish, ``How the New Mexico Draw Works,''
  https://wildlife.dgf.nm.gov/hunting/applications-and-draw-information/how-new-mexico-draw-works/.}
A hunter passed over for ten years enters the eleventh draw with the
chances of a first-time entrant.

Many institutions side with Nevada, though they disagree about how.
Arizona adds one random number per unsuccessful year; Colorado gives
most licences to whoever has the most preference points, chance only
breaking ties. Others forget, among them the lotteries that allocate
some 85,000 H-1B visas and up to 55,000 Diversity Visas a year. A worker
not selected in three successive years enters the fourth draw as a
newcomer, even under the rule in force since February 2026, which
weights H-1B registrations by the wage of the job offered.\footnote{Department
  of Homeland Security, ``Weighted Selection Process for Registrants and
  Petitioners Seeking To File Cap-Subject H-1B Petitions,'' 90
  \emph{Federal Register} 60864 (2025).}

The theory of fair allocation has not connected this disagreement to the
theory of claims, and where it touches it, it sides with New Mexico. The
best-developed account, John Broome's (\citeproc{ref-Broome1984}{1984a},
\citeproc{ref-Broome1990}{1991a}, \citeproc{ref-Broome1991}{1991b}),
holds that fairness requires claims to be satisfied in proportion to
their strength, and that when the good is indivisible and claims are
equal, an equal lottery is the fairest available procedure. It evaluates
a lottery on its own, and on its own each year's New Mexico draw is
impeccable. Broome's critics, who locate the fairness of lotteries in
what they express or exclude rather than in the chances they distribute
(\citeproc{ref-Wasserman1996}{Wasserman 1996};
\citeproc{ref-Stone2007}{Stone 2007}, \citeproc{ref-Stone2011}{2011};
\citeproc{ref-Spiekermann2022}{Spiekermann 2022}), reach the same
verdict: an impartial procedure is fair whatever last year's did. The
wider literature takes the single lottery as its unit
(\citeproc{ref-Sher1980}{Sher 1980};
\citeproc{ref-KornhauserSager1988}{Kornhauser and Sager 1988};
\citeproc{ref-Elster1989}{Elster 1989};
\citeproc{ref-Duxbury1999}{Duxbury 1999};
\citeproc{ref-Saunders2008}{Saunders 2008};
\citeproc{ref-Stone2008}{Stone 2008}). The nearest things to my thesis
are two remarks by Hersch and Rowe: that if goods are indexed to times
``we lose something that matters for fairness''
(\citeproc{ref-HerschRowe2024}{2024, sec. 5}), and that for recurring
chores ``who completed the chore last time matters to who has a claim to
avoid it this time'' (\citeproc{ref-HerschRowe2025}{2025, 3}).

I will argue for three claims. \emph{Non-separability}: the fairness of
a sequence of allocations is not a function of the fairness of its
members as the standard theory assesses them. \emph{Memory}: when units
of a good are allocated repeatedly among claimants who remain in the
pool after receiving one, fairness requires that each claimant's chances
in later rounds depend on her history of unsatisfied claims.
\emph{Currency and weight}: what memory tracks is a claimant's shortfall
from her fair share, in units of the good where a later unit is as good
to her as an earlier one and in waiting where it is not; and it may
press that shortfall only as hard as leaves no claimant's expected
share, over whatever stay she turns out to have, below what her claim
warrants. Where the good recurs, rules that forget fail the first
requirement and rules that let history decide fail the second; where
winners leave and some leave unserved, the constraint selects the rule
that forgets.

\section{Claims, chances, and division}\label{sec-broome}

A claim, in Broome's sense, is a duty owed to the candidate herself that
she should have the good, as distinct from other reasons for giving it
to her, such as that doing so would be useful
(\citeproc{ref-Broome1990}{Broome 1991a, 90--92}). Fairness requires
that ``equal claims require equal satisfaction,'' that ``stronger claims
require more satisfaction than weaker ones,'' and that ``weaker claims
must not simply be overridden by stronger ones''
(\citeproc{ref-Broome1990}{Broome 1991a, 95}).

Two features of the theory carry my argument. The first I will call the
\emph{Division Principle}. Where a good is divisible, proportional
satisfaction is achievable and required: if two people have equal claims
to a quantity of food, the food ``should be divided,'' even when giving
it all to one of them would do more good
(\citeproc{ref-Broome1990}{Broome 1991a, 95}). Broome states it most
generally in a passage that already divides a good across time:

\begin{quote}
The good, for instance, of avoiding military service could be divided by
conscripting everyone for a short time rather than a few people for a
long time. If people's claims to a good are equal, then fairness
requires that the good should be divided between them if it can be. That
is the only way to treat them with perfect equality.
(\citeproc{ref-Broome1984}{Broome 1984a, 48})
\end{quote}

The second is the account of lotteries as a second best. When the good
is indivisible and scarce, a lottery gives ``a sort of partial equality
in satisfaction'': ``each person can be given a sort of surrogate
satisfaction. By holding a lottery, each can be given an equal chance of
getting the good. This is not a perfect fairness, but it meets the
requirement of fairness to some extent''
(\citeproc{ref-Broome1990}{Broome 1991a, 97--98}). Two things follow.
Where division is available, it is required. Where it is not, the
surrogate is not merely permitted but owed: a claimant who cannot be
given her portion must be given her chance.

What duty does Nevada owe anyone to let her shoot a sheep? Broome
declines to say which reasons are claims, beyond granting that some are
(\citeproc{ref-Broome1990}{Broome 1991a, 92--93}). The claims in my
cases arise from the allocator's undertaking: a body that announces it
will allocate a public resource fairly among eligible applicants thereby
owes each of them fair treatment, and what fair treatment requires is
not the agency's to stipulate. New Mexico's announcement that it grants
no preference to past losers describes its procedure; it does not settle
whether the procedure is fair, any more than announcing that a cake will
be raffled slice by slice would. Nor does the undertaking turn every
reason for allocating a visa into a claim: the wage of an H-1B job is a
reason of usefulness, and the 2026 reform made the lottery more
responsive to usefulness, not to claims.

\section{The divisibility argument}\label{sec-division}

Consider a good that comes in indivisible units, one per period,
allocated by one agency among claimants with equal claims, where
receiving a unit does not remove a claimant from the pool: a tag a year,
a night off the rota a week. Suppose for now that the same \(N\)
claimants are present for \(T\) periods. Within any one period the
standard verdict is correct: one unit, \(N\) equal claims, no way to
divide, so hold an equal lottery. But the units are indivisible and the
sequence is not. Over \(T\) periods it is possible to give each claimant
\(T/N\) units, or as near as the integers allow. Call the procedure that
does so \emph{rotation}: draw a random order once and cycle through it.
Rotation is, in the sequence, what dividing the food is in the single
case. The Division Principle says that ``if people's claims to a good
are equal, then fairness requires that the good should be divided
between them if it can be.'' The sequence can be divided. So fairness
requires rotation, or whatever comes as close to it as circumstances
allow.

Compare New Mexico's procedure, a fresh equal lottery every period: the
\emph{forgetting} rule. Over \(T\) periods it does not divide the
sequence; it raffles it. With ten claimants and ten periods it leaves,
on average, three and a half of the ten with nothing while two or three
others receive two or more units.\footnote{The expected number with
  nothing is \(N(1-1/N)^T\). The simulation code behind every figure in
  this paper accompanies it.} Rotation leaves nobody with nothing. The
forgetting rule does over the sequence exactly what the theory forbids
in the single case: it takes a divisible good and gives it out by lot.

As a reductio: suppose someone proposed to allocate a cake among ten
people with equal claims by cutting it into ten slices and raffling each
slice separately. Each raffle, taken by itself, is what the theory
recommends for an indivisible good. Nobody thinks this is a fair way to
allocate a cake; the theory itself says the cake should be divided, and
dividing it means giving everyone a slice, not a chance at each slice.
The only difference between the cake and the tags is that the slices are
allocated at once and the tags a year apart; a cake handed out over ten
days is still a divisible cake.

Two consequences follow. The first is non-separability. The standard
theory assesses a lottery by the claims present when it is held, taking
them to be fixed by need, desert and the like, not by what earlier
lotteries did. Call a rule \emph{locally fair} at a period if its
procedure at that period, so assessed, is fair. The forgetting rule is
locally fair at every period; rotation at none after the first, since
given the history one claimant is certain to receive the unit. Yet
rotation is the fair allocation of the sequence and the forgetting rule
is not. So fairness, so assessed, is not separable across occasions. The
same result can be put another way, and from
Section~\ref{sec-extinction} on I will put it so: if claims carry
history, so that a claimant who has received less than her share has a
stronger claim to the next unit, then rotation \emph{is} locally fair at
every period and fairness is separable after all. Non-separability with
respect to history-blind claims and memory in claims are one thesis in
two vocabularies; what is not available is history-blind claims
\emph{and} separable fairness.

The second consequence is memory. Any rule that divides the sequence
must give units to those who have received fewer, which requires knowing
who has received fewer.

The argument does assume something a careful opponent will press. The
Division Principle applies to a good that ``could be divided,'' and the
tags a hunter draws in different years are not portions of one tag. What
licenses treating them so is \emph{Fungibility}: units allocated in
different periods are portions of one divisible good, for a claimant, to
the extent that her claim is to a unit of that kind and a later unit
would satisfy it as an earlier one would. The hunter's claim is to a
sheep tag, not to the 2027 tag. Fungibility is a matter of degree, since
delay is nearly always some cost, and it fails outright when a later
unit is much worse for the claimant: the housing applicant who will be
evicted first, the visa applicant whose sponsor will not wait. Then the
claim is largely a claim to timeliness. But timeliness is divisible:
delay can be shared out as units cannot, so the argument changes
currency. Fairness requires that the \emph{waiting} be divided in
proportion to claims, where the constraint of
Section~\ref{sec-proportion} allows it
(Section~\ref{sec-bottleneck}).\footnote{With unequal claims the problem
  is apportionment, and the divisor methods that solve it have a
  sequential form, which is a rule with memory
  (\citeproc{ref-WinteinHeilmann2018}{Wintein and Heilmann 2018};
  \citeproc{ref-BalinskiYoung2001}{Balinski and Young 2001}).
  ``Repeated'' fair division in computer science
  (\citeproc{ref-Igarashi2024}{Igarashi et al. 2024};
  \citeproc{ref-AleksandrovWalsh2020}{Aleksandrov and Walsh 2020}) rests
  on a related observation about envy-freeness over many rounds; it
  concerns envy, not claims.}

\section{The extinction objection}\label{sec-extinction}

The argument seems to run into a doctrine Broome states in the same
year. Once a fair lottery has given Michael the dialysis machine, it
would be absurd to hold another on the ground that lotteries are fair:
``once Michael has won the first lottery he has a perfect claim to the
treatment, and Maggie has none. Consequently it is not at all fair to
hold a second lottery'' (\citeproc{ref-Broome1984b}{Broome 1984b,
628--29}). Call this the \emph{extinction thesis}: losing a fair lottery
extinguishes the loser's claim. Vong (\citeproc{ref-Vong2015}{2015})
keeps the thesis in one form: a \emph{procedural} claim, once satisfied
by a fair lottery, is satisfied, while a \emph{benefit} claim ``cannot
be lost due to the results of a procedurally fair lottery''
(pp.~479--480); so when new units arrive the losers are ``a new case of
scarce benefit distribution'' and get a fresh fair lottery (p.~483). On
either form the persistent loser is owed nothing but a fresh lottery
among all who remain, which is the forgetting rule; and the one lottery
with memory in the literature
(\citeproc{ref-FeldblyumLeBlevennec2023}{Feldblyum Le Blevennec 2023})
protects the \emph{winner's} standing when new claimants arrive.

The reply distinguishes three structures the objection runs together. A
\emph{re-run}: the same claimants, the same unit, a second lottery
because the first went the wrong way for someone. Broome is right that
this is absurd, and the reason is that there is nothing left to divide.
An \emph{interrupted allocation}, Vong's case: one unit, and new
claimants arriving before it is handed over; here too there is one unit.
And \emph{recurrence}: a new unit each period, among overlapping
claimants. Here the extinction thesis is true and irrelevant. Maggie's
claim to the 2026 tag was extinguished when she lost; it is someone
else's. But there is a 2027 tag, and the Division Principle says the two
together are a divisible good, so the 2027 tag goes, or is more likely
to go, to someone who did not get the 2026 one. Maggie's claim to the
2027 tag is not the ghost of her claim to the 2026 tag. It is a claim to
a unit that exists, and it has more force than a past winner's for the
same reason that, when a cake is being divided, someone who has had no
slice has more of her claim outstanding than someone who has had one.
Where the currency is waiting and everyone will be served, the same
holds with waiting in place of slices.

So the extinction thesis and the memory thesis are about different
goods. To bite, Broome's objection needs a stronger thesis: that losing
a fair lottery for one unit extinguishes any bearing the loss might have
on later units. That is what the cake refutes. Someone who raffled each
slice could say of each raffle that the losers' claims to that slice had
been fairly extinguished, and it would still be true that the cake had
not been divided. Vong's form fails for a related reason: a procedural
claim is satisfied by a fresh equal lottery only if a fresh equal
lottery is the fair procedure, and over recurring units the Division
Principle says it is not. Vong rejects views on which losing weakens a
claim and on which winning strengthens it (pp.~477--479); the option he
does not consider is that losing leaves more of the claim outstanding
for \emph{later} units, and his objections, which turn on new claimants
for a single unit, do not touch it. (For his own cases, dialysis and
organs, where winners leave the pool, his verdict turns out to be right;
see Section~\ref{sec-bottleneck}.)

A residual worry makes the surrogate character of chances do real work.
If the 2026 lottery gave Maggie surrogate satisfaction, is extra weight
in 2027 not satisfying her twice? A chance is not a portion of the good;
it is what one gives when a portion cannot be given. If an equal chance
of a slice were a kind of equal slice, the raffle of the cake would be a
division of the cake, and the theory could not say, as it does, that the
cake should be divided instead. What memory counts is units received,
not chances. Wasserman's point that claimants cannot eat chances
(\citeproc{ref-Wasserman1996}{1996}) is, on this view, the beginning of
an argument for memory rather than an objection to lotteries.

\section{What memory tracks}\label{sec-proportion}

Sort the institutions by structure. In \emph{recurrence} a claimant who
receives a unit stays in the pool: tags, permits, chores. In a
\emph{bottleneck} she leaves: a visa, a flat, a kidney. Structure is
fixed by whether winners leave and currency by Fungibility, and in
practice the two run together: the claimant who leaves on being served
was waiting for one unit and cared when it came.

\subsection{Shortfall in units}\label{shortfall-in-units}

In recurrence the argument of Section~\ref{sec-division} fixes what
memory tracks: the difference between the share of the units a claimant
should have received by now and the number she has received. Call this
her \emph{shortfall}. With equal claims and one unit per period, a
claimant present for \(r\) periods who has received \(w\) units has
shortfall \(r/N - w\). Shortfall is neither the number of losses nor
time waited: points-for-losses systems reset on a win, so a claimant who
has won once in twenty years is treated as a newcomer, and they count a
loss among two hundred applicants as a loss among two; time waited
tracks shortfall only if the numbers of units and claimants are
constant.

Consider the rules that press shortfall with a weight \(\kappa \geq 0\):
each claimant's chance is proportional to
\(\max(0,\, 1 + \kappa \cdot \text{shortfall})\). At \(\kappa = 0\) this
is the forgetting rule. At \(\kappa = 1\) a claimant's expected
shortfall shrinks by a factor of \(1 - 1/N\) each period, so it is
corrected over roughly the time the practice takes to hand out one round
of shares. A past winner's weight falls below a newcomer's, and reaches
zero once her surplus reaches \(1/\kappa\); she is then excluded until
the others catch up, but never taxed, which is the extinction thesis
built in. As \(\kappa\) grows the rule approaches giving each unit to
whoever is furthest behind. In a fixed population everyone starts level,
so by symmetry each has the same chance of every unit whatever
\(\kappa\) is, and dispersion falls as \(\kappa\) rises. With ten
claimants and ten periods the expected number left with nothing falls
from 3.5 under the forgetting rule to 2.1 at \(\kappa = 1\), 1.3 at
\(\kappa = 2\) and 0 at \(\kappa = 10\). In a fixed population the
Division Principle says: press as hard as you can, and rotate.

\subsection{Newcomers and the surrogate
constraint}\label{newcomers-and-the-surrogate-constraint}

Once claimants come and go, a real objection to memory appears. In the
year she first applies, a newcomer has a claim exactly as strong as
anyone else's. Under any rule with memory she is given a smaller chance
than an incumbent who has lost before. So the rule gives unequal chances
to equal claims, which is what the standard theory says unfairness
\emph{is}.

The answer also fixes how heavily memory may press. Division of the
sequence is available only to a claimant who stays long enough for it to
be completed. A claimant who leaves after three periods cannot be given
three-tenths of a tag; for her the sequence is as indivisible as the
unit. If the allocator knew who would leave, it would owe such claimants
the surrogate outright: a chance, each period, proportional to their
claim. Not knowing, it owes the surrogate to everyone, in the only form
in which it can be given without knowing tenures in advance: that no
claimant's expected share, over whatever tenure she turns out to have,
fall short of proportionality to her claims. Call this the
\emph{surrogate constraint}. It is a floor, not a consideration to be
weighed against the dispersion that memory reduces, because the two
stand differently to the claimants. The chance is what a claimant who
will not stay for the division is \emph{owed}: for her there is nothing
to divide, the chance is all fairness can give, and a rule that gives
her less partly overrides her claim. Dispersion of shortfall among those
who stay is unfairness that fairness \emph{tolerates} because the units
cannot be divided within a period, and reducing it is an improvement in
degree. A rule may not buy an improvement of the second kind with a harm
of the first; and since the floor is owed to each claimant who may
leave, not to a class, how many such claimants there are does not bear
on whether it is owed. Within the floor, those who stay are owed as much
division as can be had. So fairness requires the rule that presses
shortfall as hard as the surrogate constraint allows.

Suppose claimants leave with a fixed probability each period,
independently of their history, with a mean tenure of ten, and the pool
is kept at ten by replacement (Table~\ref{tbl-churn}).

\begin{longtable}[]{@{}
  >{\raggedright\arraybackslash}p{(\linewidth - 8\tabcolsep) * \real{0.3939}}
  >{\raggedright\arraybackslash}p{(\linewidth - 8\tabcolsep) * \real{0.1515}}
  >{\raggedright\arraybackslash}p{(\linewidth - 8\tabcolsep) * \real{0.1515}}
  >{\raggedright\arraybackslash}p{(\linewidth - 8\tabcolsep) * \real{0.1212}}
  >{\raggedright\arraybackslash}p{(\linewidth - 8\tabcolsep) * \real{0.1818}}@{}}
\caption{Ten claimants with equal claims, one unit per period, each
leaving with probability 0.1 per period and replaced; 200 runs of 1,500
periods. A newcomer's chance in her first period (with
history-independent exit, also the expected share of claimants whose
tenure turns out to be one period) and expected shares at longer
tenures, as multiples of the proportional value \(1/N\) per period;
standard errors are below .005 in the first column and at most .02
elsewhere. Dispersion is the standard deviation, in units, of a
claimant's shortfall when she leaves.}\label{tbl-churn}\tabularnewline
\toprule\noalign{}
\begin{minipage}[b]{\linewidth}\raggedright
Rule
\end{minipage} & \begin{minipage}[b]{\linewidth}\raggedright
Newcomer's chance
\end{minipage} & \begin{minipage}[b]{\linewidth}\raggedright
Expected share, tenure 10
\end{minipage} & \begin{minipage}[b]{\linewidth}\raggedright
Tenure 30
\end{minipage} & \begin{minipage}[b]{\linewidth}\raggedright
Dispersion of realized shortfall
\end{minipage} \\
\midrule\noalign{}
\endfirsthead
\toprule\noalign{}
\begin{minipage}[b]{\linewidth}\raggedright
Rule
\end{minipage} & \begin{minipage}[b]{\linewidth}\raggedright
Newcomer's chance
\end{minipage} & \begin{minipage}[b]{\linewidth}\raggedright
Expected share, tenure 10
\end{minipage} & \begin{minipage}[b]{\linewidth}\raggedright
Tenure 30
\end{minipage} & \begin{minipage}[b]{\linewidth}\raggedright
Dispersion of realized shortfall
\end{minipage} \\
\midrule\noalign{}
\endhead
\bottomrule\noalign{}
\endlastfoot
Forgetting (\(\kappa = 0\)) & 1.00 & 0.99 & 0.97 & 0.95 \\
Shortfall, \(\kappa = 1\) & 1.01 & 1.00 & 1.01 & 0.59 \\
Shortfall, \(\kappa = 2\) & 0.97 & 0.98 & 1.00 & 0.49 \\
Shortfall, \(\kappa = 10\) & 0.56 & 0.94 & 0.98 & 0.41 \\
Points for losses, linear (Arizona) & 0.26 & 0.95 & 1.17 & 0.70 \\
Points for losses, squared (Nevada) & 0.08 & 0.99 & 1.20 & 0.58 \\
Queue / preference points (Colorado) & 0.00 & 1.01 & 1.26 & 0.47 \\
\end{longtable}

Under the forgetting rule expected share is proportional at every
tenure, 1.00 by construction (the cells are simulated); its only defect
is the dispersion of realized shortfalls, which is largest. (Dispersion
of shortfall is the right summary of the residual unfairness, since
shortfall is what the Division Principle measures; the number left with
nothing is its tail.) At \(\kappa = 1\) the newcomer's chance is
proportional to within a per cent, because a past winner's weight
reaches zero in only about one claimant-period in twenty; expected share
at every tenure is proportional to within the simulation's noise; and
dispersion falls by two-fifths. At \(\kappa = 2\) the newcomer's chance
is 97 per cent of what her claim warrants, a small but real breach; at
\(\kappa = 10\) it is little more than half.

The rules in use breach the constraint outright. Linear points give a
newcomer a quarter of the chance her claim warrants and a claimant who
stays thirty periods 17 per cent more than her share; squared points cut
the newcomer to a twelfth; a queue gives her nothing, and anyone who
leaves within a few periods leaves with nothing. These rules reduce
dispersion, some by more than the moderate shortfall rules, but by
transferring expected share from claimants whose tenure turns out short
to those whose tenure turns out long, which is to deny the short-stayers
the surrogate they are owed.

So, in recurrence, the third claim is established. Memory is required:
the forgetting rule is dominated at no cost in anyone's expected share.
But to remember lexically is to let an incumbent's accumulated claim
override a newcomer's present one. When the pool turns over, fairness
requires a rule that raises chances with shortfall at a rate that
corrects it over roughly the period in which the practice would
ordinarily deliver a share; the region around \(\kappa = 1\) has this
property, and the rules in use lie well outside it. Such a rule removes
some two-fifths of the dispersion the forgetting rule produces, and a
residual lottery remains; the residual is exactly the surrogate the
short-stayers are owed. Only where the pool is fixed is it zero, and
there the rule is rotation.

\subsection{Bottlenecks and waiting}\label{sec-bottleneck}

In a bottleneck nobody accumulates more than one unit; what varies is
time waited. Fungibility fails, because each claimant's claim is to one
unit and to having it soon; what is divisible is the waiting. Under the
forgetting rule, waiting is raffled: with a pool of ten kept full by
arrivals and one unit per period, some are served at once and about one
in seventy waits forty periods or more, though everyone's expected wait
is nine. A queue serves everyone at nine. Time, unlike a flat, is
divisible, so the Division Principle applies to the waiting and says:
equalize it. This is why John and Millum
(\citeproc{ref-JohnMillum2020}{2020}) are right that a queue can be
fairer than a lottery, and why Hersch and Rowe
(\citeproc{ref-HerschRowe2024}{2024}) are right to separate
``bottleneck'' cases, where queues belong, from one-shot ``scarcity''
cases, where lotteries do; between them lies recurrence, where a queue
gives a short-stayer none of her accruing share.

The surrogate constraint carries over, and here it cuts the other way.
Where claimants may leave unserved, by dying or losing eligibility, and
the allocator cannot tell who, those to whom the division is not
available are owed the surrogate: an equal chance, each period, at the
unit. A pure queue denies it outright. But so, to a degree, does any
rule that weights by time waited. In recurrence the newcomer's floor
could be met because past winners hold a surplus and can be pushed below
her; in a bottleneck nobody holds a surplus, since winners leave, so
every unit of chance given to a long waiter is taken from someone whose
claim is still wholly outstanding, the newcomer among them. With exits
unforeseeable and uncorrelated with waiting, the only rule under which
no tenure class's expected treatment falls below the equal chance is the
forgetting rule. So the result is two-sided in a second way. Where
everyone will be served, the waiting is divisible and fairness requires
a queue; where some will leave unserved and the allocator cannot tell
who, fairness requires the lottery that forgets. Memory is required in
the one structure and forbidden in the other, and the difference is
whether there are surplus-holders to pay for it. Housing lists and
transplant lists, where some always leave unserved, fall on the second
side: a queue is fairer than a lottery, as John and Millum say, exactly
to the extent that exit is foreseeable. (Where an exit is foreseeable
and imminent the claim is urgent and the rule must be weighted for it,
as transplant allocation is; that is a difference in claims, not a
memory of history.)\footnote{Dynamic-allocation economics reaches
  related conclusions by another route (\citeproc{ref-Leshno2022}{Leshno
  2022}; \citeproc{ref-ArnostiShi2020}{Arnosti and Shi 2020}), but on
  grounds of efficiency rather than of what applicants are owed.}

The verdict for the visa lotteries follows, and it is not the one the
opening cases suggest. They are bottlenecks with attrition: a selected
worker leaves the pool, and an unselected one may lose her sponsor or
her status in ways the allocator cannot foresee. So the lottery that
forgets is the fair rule, and Hersch and Rowe's
(\citeproc{ref-HerschRowe2024}{2024, secs. 1, 3}) endorsement of it is
right, though not for their reason, that the case is scarcity rather
than bottleneck; it is right because the surrogate owed to those who
will leave unserved cannot be paid by anyone else. What is wrong with
the rule in force since 2026 is not that it forgets but that it weights
by the wage of the job, a reason of usefulness rather than a claim. And
their remark that a periodic lottery in a bottleneck is unfair because a
person ``who arrives into the lottery pool early and is continuously
unlucky'' may wait far longer than a lucky newcomer (sec.~4) is true
only where everyone will be served.

\section{How far does memory reach?}\label{sec-reach}

If Nevada must remember Maggie's 2026 loss in 2027, must it remember her
loss in the housing lottery? In Arizona's draw? Her mother's? A view
that cannot stop somewhere principled collapses into a demand that luck
be equalized, a different and far more contested thesis. Three bounds,
besides the bound in weight that Section~\ref{sec-proportion}
established, follow from the argument that generates the claim.

\emph{Memory is bound to the practice.} The corrective claim arises
because a particular agency's allocation of a particular good left
residual unfairness; it is a claim on the division of \emph{those}
units, addressed to \emph{that} agency. The good matters because the
2027 tag can divide the sequence of tags only because it is a tag; a
flat is not a portion of what the hunter has a claim to. Goods are
individuated by claims, not names: a Nevada tag and an Arizona tag are
different goods, because a claim to hunt in one state is not satisfied
by a tag for the other, whereas two draws run by one agency for the same
tag are one good, so an allocator cannot shed its obligation by
splitting a practice in two. The agency matters because a division is
something an agent does with what it has; another allocator has none of
its units to give. This distinguishes the thesis from luck
egalitarianism (\citeproc{ref-Dworkin1981}{Dworkin 1981};
\citeproc{ref-Cohen1989}{Cohen 1989}; \citeproc{ref-Temkin1993}{Temkin
1993}), whose currency is compensation, so that a shortfall in one good
can be made up in another; mine concerns the fair division of one good
by the agency dividing it, in that good's currency, and says nothing
about whether the persistent loser is worse off overall.

\emph{Memory is bound to the claim.} The corrective claim is a
strengthened claim to the good, not a credit. Someone who applied
unsuccessfully for ten years and then gave up hunting is owed nothing;
there is no tag she has a claim to. Memory is not compensation, since
the lottery did no wrong. Two boundary questions have determinate
answers. \emph{Transition}: when a forgetting practice adopts memory,
losses under the old rule count, since the shortfall existed whether or
not the rule tracked it; a rectificatory view would say the opposite.
\emph{Regress}: the 2028 allocation need not correct the 2027
allocation's failure to correct fully, because shortfall is a state, not
a stack of debts; at each period there is one quantity, and correcting
it is the whole of what is owed.

\emph{Memory is bound to the claimant.} The claimant is by default the
person, because claims are owed to someone; New Mexico cannot make its
persistent losers vanish by declaring that its claimants are ``this
year's applicants.'' Nor is there an obligation to correct one's
ancestors' shortfalls, which would be claims of rectification; the
lottery did no wrong.

\section{Objections}\label{sec-objections}

\emph{Losing on purpose.} The objection that memory rewards losing is
familiar: the NBA's draft lottery rewarded losing and was flattened for
that reason (cf. \citeproc{ref-Elster1988}{Elster 1988}). In the cases
that motivate this paper, losing is not under the claimant's control:
she cannot make herself lose, only apply, and applying without a genuine
claim is the problem of verifying claims, which every allocation faces.
Where losing is controllable, as in a league, the objection shows that
fairness is in tension with incentive, which is a reason to weigh
fairness against other things, not evidence that fairness requires
forgetting. What must be conceded is that memory raises the value of a
manufactured claim: a speculative entrant gains a ticket that grows. The
response is to attach memory to participation that costs something, as a
practice does when it requires a paid licence to apply. That is a
constraint on implementing fairness, not evidence that fairness has no
memory.

\emph{Equality, not fairness.} Someone might grant the verdict on the
outcomes and deny that fairness is what condemns them: what is bad about
the forgetting rule is inequality, or that some are badly off, matters
for a theory of equality or priority (\citeproc{ref-Parfit2000}{Parfit
2000}). Suppose the persistent loser is, all things considered, no worse
off than the repeat winner. Her complaint survives: this agency,
dividing this good among people with equal claims, has given her less
than her share. That complaint is comparative, concerns the good in
question, and is addressed to the allocator, the marks of a fairness
complaint in Broome's sense; a priority view could not explain why it
lapses when she stops claiming the good.

\emph{Ex ante.} The last objection comes from the view, descending from
Diamond's (\citeproc{ref-Diamond1967}{1967}) reply to Harsanyi
(\citeproc{ref-Harsanyi1955}{1955}) (\citeproc{ref-Frick2015}{Frick
2015}; \citeproc{ref-OtsukaVoorhoeve2009}{Otsuka and Voorhoeve 2009}),
that the fairness of a procedure is a matter of the prospects it gives
people at the time of choice. On that view the forgetting rule and
rotation are equally fair, since at the outset each claimant has the
same prospects under both; the forgetting rule's defect, as the first
row of Table~\ref{tbl-churn} showed, is entirely ex post. First reply:
the ex ante view must say when the ex ante is. If it is the beginning of
the practice, the view cannot prefer the forgetting rule to rotation,
since they are ex ante identical, and the Division Principle settles
what it leaves open. If it is the beginning of each period, the view is
asserting separability with respect to history-blind claims, which is
what the cake refutes. Second, the ex ante view was developed for
one-shot cases, where an outcome cannot be undone and prospects are all
a procedure can distribute. In recurrence what was a prospect becomes a
fact the next allocation can attend to.

\section{Conclusion}\label{sec-conclusion}

Fairness has memory where the good recurs. What the memory tracks is a
claimant's shortfall from her share of the good being divided; it is
owed by the agency dividing the good, for as long as she has a claim to
it; and it may be pressed only as hard as leaves no claimant's expected
share, over whatever stay she turns out to have, below what her claim
warrants. Where winners leave and some will leave unserved, the same
constraint forbids memory, because there is no one whose surplus can pay
for it. The result was visible in the example Broome chose to illustrate
the Division Principle: military service, divided by conscripting
everyone for a short time. Everyone's short time is a rotation. The
lottery is what one does when one cannot rotate, and one can rotate more
often than the lottery's defenders have noticed.

\protect\phantomsection\label{refs}
\begin{CSLReferences}{1}{0}
\bibitem[\citeproctext]{ref-AleksandrovWalsh2020}
Aleksandrov, Martin, and Toby Walsh. 2020. {``Online Fair Division: A
Survey.''} In \emph{Proceedings of the AAAI Conference on Artificial
Intelligence}, 34:13557--62. 9. doi:
\href{https://doi.org/10.1609/aaai.v34i09.7081}{10.1609/aaai.v34i09.7081}.

\bibitem[\citeproctext]{ref-ArnostiShi2020}
Arnosti, Nick, and Peng Shi. 2020. {``Design of Lotteries and Wait-Lists
for Affordable Housing Allocation.''} \emph{Management Science} 66 (6):
2291--2307. doi:
\href{https://doi.org/10.1287/mnsc.2019.3311}{10.1287/mnsc.2019.3311}.

\bibitem[\citeproctext]{ref-BalinskiYoung2001}
Balinski, Michel L., and H. Peyton Young. 2001. \emph{Fair
Representation: Meeting the Ideal of One Man, One Vote}. 2nd ed.
Washington, DC: Brookings Institution Press. First edition Yale
University Press, 1982.

\bibitem[\citeproctext]{ref-Broome1984}
Broome, John. 1984a. {``Selecting People Randomly.''} \emph{Ethics} 95
(1): 38--55. doi: \href{https://doi.org/10.1086/292596}{10.1086/292596}.

\bibitem[\citeproctext]{ref-Broome1984b}
---------. 1984b. {``Uncertainty and Fairness.''} \emph{The Economic
Journal} 94 (375): 624--32. doi:
\href{https://doi.org/10.2307/2232707}{10.2307/2232707}.

\bibitem[\citeproctext]{ref-Broome1990}
---------. 1991a. {``Fairness.''} \emph{Proceedings of the Aristotelian
Society} 91: 87--101. doi:
\href{https://doi.org/10.1093/aristotelian/91.1.87}{10.1093/aristotelian/91.1.87}.

\bibitem[\citeproctext]{ref-Broome1991}
---------. 1991b. \emph{Weighing Goods: Equality, Uncertainty and Time}.
Oxford: Blackwell.

\bibitem[\citeproctext]{ref-Cohen1989}
Cohen, G. A. 1989. {``On the Currency of Egalitarian Justice.''}
\emph{Ethics} 99 (4): 906--44. doi:
\href{https://doi.org/10.1086/293126}{10.1086/293126}.

\bibitem[\citeproctext]{ref-Diamond1967}
Diamond, Peter A. 1967. {``Cardinal Welfare, Individualistic Ethics, and
Interpersonal Comparison of Utility: Comment.''} \emph{Journal of
Political Economy} 75 (5): 765--66. doi:
\href{https://doi.org/10.1086/259353}{10.1086/259353}.

\bibitem[\citeproctext]{ref-Duxbury1999}
Duxbury, Neil. 1999. \emph{Random Justice: On Lotteries and Legal
Decision-Making}. Oxford: Oxford University Press.

\bibitem[\citeproctext]{ref-Dworkin1981}
Dworkin, Ronald. 1981. {``What Is Equality? Part 2: Equality of
Resources.''} \emph{Philosophy \& Public Affairs} 10 (4): 283--345.

\bibitem[\citeproctext]{ref-Elster1988}
Elster, Jon. 1988. {``Taming Chance: Randomization in Individual and
Social Decisions.''} In \emph{The Tanner Lectures on Human Values},
edited by Sterling M. McMurrin, 9:105--80. Salt Lake City: University of
Utah Press.

\bibitem[\citeproctext]{ref-Elster1989}
---------. 1989. \emph{Solomonic Judgements: Studies in the Limitations
of Rationality}. Cambridge: Cambridge University Press.

\bibitem[\citeproctext]{ref-FeldblyumLeBlevennec2023}
Feldblyum Le Blevennec, Marie Kerguelen. 2023. {``Fairness and Chance in
Diachronic Lotteries: A Response to Vong.''} \emph{Journal of Ethics and
Social Philosophy} 26 (1): 192--203. doi:
\href{https://doi.org/10.26556/jesp.v26i1.2190}{10.26556/jesp.v26i1.2190}.

\bibitem[\citeproctext]{ref-Frick2015}
Frick, Johann. 2015. {``Contractualism and Social Risk.''}
\emph{Philosophy \& Public Affairs} 43 (3): 175--223. doi:
\href{https://doi.org/10.1111/papa.12058}{10.1111/papa.12058}.

\bibitem[\citeproctext]{ref-Harsanyi1955}
Harsanyi, John C. 1955. {``Cardinal Welfare, Individualistic Ethics, and
Interpersonal Comparisons of Utility.''} \emph{Journal of Political
Economy} 63 (4): 309--21. doi:
\href{https://doi.org/10.1086/257678}{10.1086/257678}.

\bibitem[\citeproctext]{ref-HerschRowe2024}
Hersch, Gil, and Thomas Rowe. 2024. {``Lotteries, Queues, and
Bottlenecks.''} In \emph{Oxford Studies in Political Philosophy}, edited
by David Sobel and Steven Wall, 10:186--210. Oxford University Press.
doi:
\href{https://doi.org/10.1093/oso/9780198909460.003.0008}{10.1093/oso/9780198909460.003.0008}.

\bibitem[\citeproctext]{ref-HerschRowe2025}
---------. 2025. {``Allocative Fairness.''} \emph{Philosophy Compass} 20
(5): e70042. doi:
\href{https://doi.org/10.1111/phc3.70042}{10.1111/phc3.70042}.

\bibitem[\citeproctext]{ref-Igarashi2024}
Igarashi, Ayumi, Martin Lackner, Oliviero Nardi, and Arianna Novaro.
2024. {``Repeated Fair Allocation of Indivisible Items.''} In
\emph{Proceedings of the AAAI Conference on Artificial Intelligence},
38:9781--89. 9. doi:
\href{https://doi.org/10.1609/aaai.v38i9.28837}{10.1609/aaai.v38i9.28837}.

\bibitem[\citeproctext]{ref-JohnMillum2020}
John, Tyler M., and Joseph Millum. 2020. {``First Come, First Served?''}
\emph{Ethics} 130 (2): 179--207. doi:
\href{https://doi.org/10.1086/705763}{10.1086/705763}.

\bibitem[\citeproctext]{ref-KornhauserSager1988}
Kornhauser, Lewis A., and Lawrence G. Sager. 1988. {``Just Lotteries.''}
\emph{Social Science Information} 27 (4): 483--516. doi:
\href{https://doi.org/10.1177/053901888027004001}{10.1177/053901888027004001}.

\bibitem[\citeproctext]{ref-Leshno2022}
Leshno, Jacob D. 2022. {``Dynamic Matching in Overloaded Waiting
Lists.''} \emph{American Economic Review} 112 (12): 3876--3910. doi:
\href{https://doi.org/10.1257/aer.20201111}{10.1257/aer.20201111}.

\bibitem[\citeproctext]{ref-OtsukaVoorhoeve2009}
Otsuka, Michael, and Alex Voorhoeve. 2009. {``Why It Matters That Some
Are Worse Off Than Others: An Argument Against the Priority View.''}
\emph{Philosophy \& Public Affairs} 37 (2): 171--99. doi:
\href{https://doi.org/10.1111/j.1088-4963.2009.01154.x}{10.1111/j.1088-4963.2009.01154.x}.

\bibitem[\citeproctext]{ref-Parfit2000}
Parfit, Derek. 2000. {``Equality or Priority?''} In \emph{The Ideal of
Equality}, edited by Matthew Clayton and Andrew Williams, 81--125.
Basingstoke: Macmillan.

\bibitem[\citeproctext]{ref-Saunders2008}
Saunders, Ben. 2008. {``The Equality of Lotteries.''} \emph{Philosophy}
83 (3): 359--72. doi:
\href{https://doi.org/10.1017/S0031819108000727}{10.1017/S0031819108000727}.

\bibitem[\citeproctext]{ref-Sher1980}
Sher, George. 1980. {``What Makes a Lottery Fair?''} \emph{No{û}s} 14
(2): 203--16. doi:
\href{https://doi.org/10.2307/2214861}{10.2307/2214861}.

\bibitem[\citeproctext]{ref-Spiekermann2022}
Spiekermann, Kai. 2022. {``Good Reasons for Losers: Lottery
Justification and Social Risk.''} \emph{Economics and Philosophy} 38
(1): 108--31. doi:
\href{https://doi.org/10.1017/S0266267121000043}{10.1017/S0266267121000043}.

\bibitem[\citeproctext]{ref-Stone2007}
Stone, Peter. 2007. {``Why Lotteries Are Just.''} \emph{Journal of
Political Philosophy} 15 (3): 276--95. doi:
\href{https://doi.org/10.1111/j.1467-9760.2006.00274.x}{10.1111/j.1467-9760.2006.00274.x}.

\bibitem[\citeproctext]{ref-Stone2008}
---------. 2008. {``On Fair Lotteries.''} \emph{Social Theory and
Practice} 34 (4): 573--90. doi:
\href{https://doi.org/10.5840/soctheorpract200834431}{10.5840/soctheorpract200834431}.

\bibitem[\citeproctext]{ref-Stone2011}
---------. 2011. \emph{The Luck of the Draw: The Role of Lotteries in
Decision Making}. New York: Oxford University Press. doi:
\href{https://doi.org/10.1093/acprof:oso/9780199756100.001.0001}{10.1093/acprof:oso/9780199756100.001.0001}.

\bibitem[\citeproctext]{ref-Temkin1993}
Temkin, Larry S. 1993. \emph{Inequality}. New York: Oxford University
Press.

\bibitem[\citeproctext]{ref-Vong2015}
Vong, Gerard. 2015. {``Fairness, Benefiting by Lottery and the Chancy
Satisfaction of Moral Claims.''} \emph{Utilitas} 27 (4): 470--86. doi:
\href{https://doi.org/10.1017/S0953820815000357}{10.1017/S0953820815000357}.

\bibitem[\citeproctext]{ref-Wasserman1996}
Wasserman, David. 1996. {``Let Them Eat Chances: Probability and
Distributive Justice.''} \emph{Economics and Philosophy} 12 (1): 29--49.
doi:
\href{https://doi.org/10.1017/S0266267100003709}{10.1017/S0266267100003709}.

\bibitem[\citeproctext]{ref-WinteinHeilmann2018}
Wintein, Stefan, and Conrad Heilmann. 2018. {``Dividing the Indivisible:
Apportionment and Philosophical Theories of Fairness.''} \emph{Politics,
Philosophy \& Economics} 17 (1): 51--74. doi:
\href{https://doi.org/10.1177/1470594X17715248}{10.1177/1470594X17715248}.

\end{CSLReferences}




\end{document}
